\documentclass[12pt,a4paper]{article}

% ── Pacotes ────────────────────────────────────────────────────────────────────
\usepackage[utf8]{inputenc}
\usepackage[T1]{fontenc}
\usepackage[english]{babel}
\usepackage{amsmath,amssymb,amsthm}
\usepackage{mathrsfs}
\usepackage{geometry}
\usepackage{xcolor}
\usepackage{hyperref}
\usepackage{booktabs}
\usepackage{array}
\usepackage{enumitem}
\usepackage{mdframed}

\geometry{
	a4paper,
	top=2.5cm, bottom=2.5cm,
	left=3cm, right=3cm
}

\hypersetup{
	colorlinks=true,
	linkcolor=blue,
	citecolor=blue,
	urlcolor=blue,
	pdftitle={Motor de Herança Estrutural — Documento de Formalização (Versão Definitiva)},
	pdfauthor={Tiago Bandeira}
}

% ── Ambientes de Teorema ───────────────────────────────────────────────────────
\newtheorem{theorem}{Teorema}[section]
\newtheorem{lemma}[theorem]{Lema}
\newtheorem{proposition}[theorem]{Proposição}
\newtheorem{corollary}[theorem]{Corolário}

\theoremstyle{definition}
\newtheorem{definition}[theorem]{Definição}
\newtheorem{example}[theorem]{Exemplo}
\newtheorem{remark}[theorem]{Observação}

% ── Atalhos matemáticos ────────────────────────────────────────────────────────
\newcommand{\N}{\mathbb{N}}
\newcommand{\Z}{\mathbb{Z}}
\newcommand{\R}{\mathbb{R}}
\newcommand{\C}{\mathbb{C}}
\newcommand{\PP}{\mathcal{P}}          % conjunto dos primos
\newcommand{\F}{\mathcal{F}}           % fita-dobra
\newcommand{\G}{\mathcal{G}}           % grade
\newcommand{\HRm}{\mathrm{HR}^{-}}
\newcommand{\HRp}{\mathrm{HR}^{+}}
\newcommand{\norm}[1]{\left\| #1 \right\|}
\newcommand{\abs}[1]{\left| #1 \right|}

% ── Título ─────────────────────────────────────────────────────────────────────
\title{%
	\textbf{Motor de Herança Estrutural}\\[0.4em]
	\large Documento de Formalização — Versão Definitiva (rev.\ 2)
}
\author{Tiago Bandeira\\[0.3em]
	\small\texttt{tiagobandeira.goldbach2025@gmail.com}
}
\date{Maio de 2026}

% ═══════════════════════════════════════════════════════════════════════════════
\begin{document}

\maketitle
\thispagestyle{empty}

% ── Nota do Autor ──────────────────────────────────────────────────────────────
\begin{center}
	\fbox{\parbox{0.88\textwidth}{%
		\textbf{Nota do Autor.} Este documento substitui e corrige as
		definições do Motor presentes nos Artigos~5 e~6 da série. A distinção
		central corrigida nesta revisão é: o par base~$2M$ é a soma das colunas
		da fita — o sistema já sabe que é Goldbach. O par alvo~$2M{+}2$ é a
		soma das diagonais da fita — quem certifica que é Goldbach é
		$\HRm$. Esta distinção propaga-se por todas as definições do motor.
	}}
\end{center}
\vspace{0.5em}

\tableofcontents
\newpage

% ═══════════════════════════════════════════════════════════════════════════════
\section*{0.\quad Glossário — Distinção Fundamental}
\addcontentsline{toc}{section}{0.\enspace Glossário — Distinção Fundamental}
\label{sec:glossario}

Antes de qualquer definição, é necessário fixar com precisão os dois
conceitos que organizam todo o motor:

\renewcommand{\arraystretch}{1.5}
\begin{center}
\begin{tabular}{>{\bfseries}p{2.8cm} p{5.8cm} p{4.6cm}}
	\toprule
	\textbf{Conceito} & \textbf{Definição} & \textbf{Quem garante} \\
	\midrule
	Par base $2M$ &
		Soma das colunas da fita: $a_k + b_k = 2M$ &
		O sistema escolhe colunas com primos — Goldbach já é conhecido para
		$2M$. \\[0.4em]
	Par alvo $2M{+}2$ &
		Soma das diagonais da fita:
		$a_{(k+1)} + b_k = 2M{+}2$,\newline $k \in \{1, \ldots, N{-}1\}$ &
		$\HRm$ certifica que $2M{+}2$ é Goldbach ao encontrar primos nos
		extremos de $W_i$. \\
	\bottomrule
\end{tabular}
\end{center}

% ═══════════════════════════════════════════════════════════════════════════════
\section{Motivação e Mudança de Paradigma}
\label{sec:motivacao}

A Conjectura de Goldbach afirma que todo número par maior que~2 é soma de
dois primos. A abordagem clássica trata essa questão de forma estática —
para cada par~$2M$, pergunta-se se existe uma decomposição $2M = p + q$.
O Motor de Herança Estrutural substitui essa pergunta estática pela
pergunta dinâmica:

\begin{center}
	\fbox{\parbox{0.85\textwidth}{%
		\itshape Dado que $2M$ é Goldbach (par base), \textbf{a estrutura
		geométrica garante que $2M{+}2$ (par alvo) também o é?}
	}}
\end{center}

\noindent O motor organiza os ímpares em dois objetos acoplados — a
Fita-Dobra e a Grade — de forma que a certificação de Goldbach para o par
alvo deixe de ser uma questão aritmética pontual e passe a ser uma
propriedade emergente de ocupação espacial. A varredura geométrica
substitui o teste aritmético.

% ═══════════════════════════════════════════════════════════════════════════════
\section{A Fita-Dobra \texorpdfstring{$\mathcal{F}(2{\times}N)$}{F(2xN)}}
\label{sec:fita}

\subsection{Definição}
\label{subsec:fita-def}

A Fita-Dobra é uma matriz $2{\times}N$ construída a partir do par base~$2M$.
A Linha~1 percorre os ímpares crescentes da esquerda para a direita. A
Linha~2 percorre os complementares da direita para a esquerda, de modo que
cada coluna~$k$ satisfaz $a_k + b_k = 2M$ para todo~$k$:

\renewcommand{\arraystretch}{1.4}
\begin{center}
\begin{tabular}{>{\bfseries}l ccccc}
	\toprule
	\textbf{Linha 1} $\rightarrow$ & $a_1 = 1$ & $a_2 = 3$ & $a_3 = 5$ & $\cdots$ & $a_N = 2N{-}1$ \\
	\midrule
	Linha 2 $\leftarrow$ & $b_1 = 2M{-}1$ & $b_2 = 2M{-}3$ & $b_3 = 2M{-}5$ & $\cdots$ & $b_N = 2M{-}(2N{-}1)$ \\
	\bottomrule
\end{tabular}
\end{center}

\subsection{Propriedade Fundamental — Soma Constante (Par Base)}
\label{subsec:fita-soma}

\begin{proposition}[Par Base]
	Para toda coluna $k \in \{1, \ldots, N\}$:
	\[
		a_k + b_k = 2M.
	\]
	Esta soma é constante, independente da posição~$k$ e independente de
	primalidade. A fita é cega à posição e cega à primalidade. O par
	base~$2M$ é certificado por Goldbach ao se encontrar uma coluna~$k$ com
	$a_k$ e $b_k$ ambos primos — o motor parte desse fato como dado.
\end{proposition}

\subsection{As Diagonais da Fita — Par Alvo}
\label{subsec:fita-diagonais}

Além dos pares colunares, a fita contém pares diagonais: pares
$(a_{(k+1)}, b_k)$ que cruzam entre colunas adjacentes. Por construção:

\begin{center}
	\fbox{\parbox{0.85\textwidth}{%
		\begin{align*}
		a_{(k+1)} + b_k
		&= \bigl(2(k{+}1){-}1\bigr) + (2M{-}2k{+}1) \\
		&= (2k{+}1) + (2M{-}2k{+}1) \;=\; 2M{+}2 \;=\; \text{par alvo}
		\end{align*}
	}}
\end{center}

\noindent Cada par diagonal soma exatamente $2M{+}2$, independente de
primalidade. Esses pares diagonais são o objeto que o scanner captura e
armazena em~$W_i$. A certificação de que $2M{+}2$ é Goldbach depende de
encontrar uma diagonal com ambos os extremos primos — isso é precisamente
$\HRm$.

\subsection{A Fita Representa Dois Pares Consecutivos}
\label{subsec:fita-dois-pares}

A Fita-Dobra representa simultaneamente dois números pares consecutivos: o
par base~$2M$ (soma das colunas) e o par alvo~$2M{+}2$ (soma das
diagonais). Isso é necessário para que o motor cubra todos os pares sem
exceção.

\begin{remark}
	Quando $M$ é ímpar, $M + M = 2M$ faz com que o centro da fita
	(coluna~$N$) tenha um elemento repetido. Quando $M$ é par, o centro tem
	$M$ e $M{+}1$ adjacentes sem repetição. Esta adaptação garante a
	proporção $2{:}1$ entre naturais e ímpares e que a fita cubra dois pares
	consecutivos distintos.
\end{remark}

% ═══════════════════════════════════════════════════════════════════════════════
\section{A Grade \texorpdfstring{$\mathcal{G}(3{\times}C)$}{G(3xC)} e o Acoplamento}
\label{sec:grade}

\subsection{Condição de Acoplamento}
\label{subsec:grade-acoplamento}

A Grade $\mathcal{G}(3{\times}C)$ é uma matriz de três linhas e $C$ colunas
que contém os mesmos ímpares da Fita-Dobra, organizados em percurso zigzag.
Para que grade e fita representem exatamente o mesmo conjunto de ímpares, a
condição de acoplamento deve ser satisfeita:

\begin{center}
	\fbox{\parbox{0.65\textwidth}{%
		\centering
		$3C = 2N$\\[0.3em]
		\small O número de células da grade ($3C$) deve igualar o número de
		células da fita ($2N$).
	}}
\end{center}

\noindent Quando o número total de ímpares não é divisível por~3,
duplicações centrais são introduzidas para satisfazer a condição sem alterar
a cardinalidade efetiva do sistema.

\subsection{Percurso Zigzag da Grade}
\label{subsec:grade-zigzag}

Os ímpares são dispostos na grade em percurso zigzag de três linhas:
Linha~1~($\rightarrow$), Linha~2~($\leftarrow$), Linha~3~($\rightarrow$).
Este percurso garante que as diagonais geométricas da grade correspondam às
diagonais da fita.

\renewcommand{\arraystretch}{1.4}
\begin{center}
\begin{tabular}{>{\bfseries}l cccccccc}
	\toprule
	Linha 1 $\rightarrow$ & 1 & 3 & 5 & 7 & 9 & 11 & 13 & 15 \\
	\midrule
	Linha 2 $\leftarrow$  & 31 & 29 & 27 & 25 & 23 & 21 & 19 & 17 \\
	\midrule
	Linha 3 $\rightarrow$ & 33 & 35 & 37 & 39 & 41 & 43 & 45 & 47 \\
	\bottomrule
\end{tabular}
\end{center}

\noindent\textit{Exemplo: Grade $\mathcal{G}(3{\times}8)$. Par base
$2M = 48$, par alvo $2M{+}2 = 50$. Acoplamento: $3{\times}8 = 2{\times}12
= 24$ células.~$\checkmark$}

% ═══════════════════════════════════════════════════════════════════════════════
\section{A Janela \texorpdfstring{$W_i$}{Wi} — Definição Correta}
\label{sec:janela}

A Janela~$W_i$ é o objeto central do motor. Ela é uma matriz $3{\times}3$
fixa que armazena, nos seus quatro eixos de simetria, até quatro diagonais
da fita que somam o par alvo~$2M{+}2$. $W_i$ não é deslizante: ela é
preenchida pelo scanner e permanece como referência estrutural.

\subsection{Os Quatro Eixos de Simetria de \texorpdfstring{$W_i$}{Wi}}
\label{subsec:janela-eixos}

Uma matriz $3{\times}3$ possui exatamente quatro linhas de simetria que
passam pelo centro. Cada eixo conecta dois extremos opostos e define um
slot para um par diagonal da fita (que soma o par alvo~$2M{+}2$):

\renewcommand{\arraystretch}{1.5}
\begin{center}
\begin{tabular}{>{\bfseries}p{3.8cm} p{8cm}}
	\toprule
	D1 — Diagonal Principal & Canto superior esquerdo $[0,0]$ $\leftrightarrow$ Canto inferior direito $[2,2]$ \\
	\midrule
	D2 — Diagonal Secundária & Canto superior direito $[0,2]$ $\leftrightarrow$ Canto inferior esquerdo $[2,0]$ \\
	\midrule
	Cv — Coluna Central & Topo central $[0,1]$ $\leftrightarrow$ Base central $[2,1]$ \\
	\midrule
	Lc — Linha Central & Esquerda central $[1,0]$ $\leftrightarrow$ Direita central $[1,2]$ \\
	\bottomrule
\end{tabular}
\end{center}

\noindent Em cada eixo: o elemento menor da diagonal é posicionado na
entrada, o maior na posição oposta. O centro $[1,1]$ armazena o pivô — o
elemento central da varredura do scanner.

\subsection{Estrutura Visual de \texorpdfstring{$W_i$}{Wi}}
\label{subsec:janela-visual}

\renewcommand{\arraystretch}{1.6}
\begin{center}
\begin{tabular}{|c|c|c|}
	\hline
	$[0,0]$\; $E_1(\mathrm{D1})$ & $[0,1]$\; $E_1(\mathrm{Cv})$ & $[0,2]$\; $E_1(\mathrm{D2})$ \\
	\hline
	$[1,0]$\; $E_1(\mathrm{Lc})$ & $[1,1]$\; \textbf{PIVÔ} & $[1,2]$\; $E_2(\mathrm{Lc})$ \\
	\hline
	$[2,0]$\; $E_2(\mathrm{D2})$ & $[2,1]$\; $E_2(\mathrm{Cv})$ & $[2,2]$\; $E_2(\mathrm{D1})$ \\
	\hline
\end{tabular}
\end{center}

\noindent\textit{Em cada eixo: $E_1 + E_2 = 2M{+}2$ (par alvo) —
garantido incondicionalmente pela construção da fita.}

\subsection{Exemplo Numérico de \texorpdfstring{$W_i$}{Wi} (par base = 48, par alvo = 50)}
\label{subsec:janela-exemplo}

O scanner percorre a grade canônica do par base~48 e preenche $W_i$ com
quatro diagonais que somam o par alvo~50. O elemento~1 é ignorado pois não
forma diagonal secundária válida.

\renewcommand{\arraystretch}{1.6}
\begin{center}
\begin{tabular}{|c|c|c|}
	\hline
	7 & 9 & 11 \\
	\hline
	37 & 25 & 13 \\
	\hline
	39 & 41 & 43 \\
	\hline
\end{tabular}
\end{center}

\noindent
$\mathrm{D1}$: $7 + 43 = 50\ \checkmark$\quad
$\mathrm{Cv}$: $9 + 41 = 50\ \checkmark$\quad
$\mathrm{D2}$: $11 + 39 = 50\ \checkmark$\quad
$\mathrm{Lc}$: $13 + 37 = 50\ \checkmark$

\medskip\noindent
\textit{$\HRm$: verificar se algum par $(E_1, E_2)$ num eixo de $W_i$ tem
ambos os elementos primos. Ex.: $\mathrm{D1} \to (7, 43)$: $7$ primo
$\checkmark$, $43$ primo $\checkmark$ $\Rightarrow$ $\HRm$ ativada,
certificando $50 = 7 + 43$.}

\subsection{Número de Janelas \texorpdfstring{$W_i$}{Wi} por Varredura}
\label{subsec:janela-numero}

Como cada $W_i$ comporta até 4 diagonais e existem $N{-}1$ diagonais
internas na fita (excluindo a posição do elemento~1), o número de janelas
necessárias para cobrir todas as diagonais é aproximadamente
$\lceil(N{-}1)/4\rceil$. Quando $N{-}1$ não é divisível por~4, a varredura
final cobre o restante parcialmente, garantindo que nenhuma diagonal fique
sem registro.

% ═══════════════════════════════════════════════════════════════════════════════
\section{O Scanner — Preenchimento de \texorpdfstring{$W_i$}{Wi}}
\label{sec:scanner}

\subsection{O Mecanismo dos Pivôs}
\label{subsec:scanner-pivos}

O scanner percorre a grade e preenche os eixos de $W_i$ com as diagonais
da fita. Dois pivôs marcham das extremidades da lista de ímpares em direção
ao centro. O pivô esquerdo pula a primeira posição (elemento~1), pois o~1
não forma diagonal secundária válida — não existe elemento na grade tal que
$1 + a = \text{par alvo}\ 2M{+}2$.

\begin{center}
	\fbox{\parbox{0.85\textwidth}{%
		\textbf{Algoritmo do Scanner:}
		\begin{enumerate}[label=\arabic*., leftmargin=1.5em, itemsep=0pt]
			\item Posicionar pivô esquerdo na segunda posição da lista (pula o
			elemento~1).
			\item Posicionar pivô direito na extremidade final da lista.
			\item Registrar o par $(E_{\mathrm{esq}}, E_{\mathrm{dir}})$ num
			eixo de $W_i$. O elemento central da lista é o pivô do sensor.
			\item Avançar pivô esquerdo ($+1$) e recuar pivô direito ($-1$).
			\item Repetir até os pivôs se encontrarem. Cada iteração preenche
			um eixo de $W_i$.
			\item Quando $W_i$ tiver os 4 eixos preenchidos, registrá-la e
			iniciar nova $W_i$ para o bloco seguinte.
		\end{enumerate}
	}}
\end{center}

\subsection{Propriedade de Soma Constante nos Eixos}
\label{subsec:scanner-soma}

\begin{proposition}[Soma Constante nos Eixos de $W_i$]
	Por construção da Fita-Dobra (Seção~\ref{subsec:fita-diagonais}), cada
	par $(E_1, E_2)$ capturado pelo scanner e armazenado num eixo de $W_i$
	satisfaz automaticamente
	\[
		E_1 + E_2 = 2M{+}2 \quad (\text{par alvo}),
	\]
	independente de os elementos serem ou não primos. Esta é a garantia
	algébrica incondicional do motor.
\end{proposition}

% ═══════════════════════════════════════════════════════════════════════════════
\section{As Hipóteses Restritivas \texorpdfstring{$\HRm$}{HR-} e \texorpdfstring{$\HRp$}{HR+}}
\label{sec:HR}

O motor opera em duas camadas. A camada algébrica — construção da fita,
acoplamento com a grade, preenchimento de $W_i$ pelo scanner — é
inteiramente incondicional. A camada aritmética introduz as únicas
hipóteses não provadas:

\renewcommand{\arraystretch}{1.5}
\begin{center}
\begin{tabular}{>{\bfseries}p{2.8cm} p{10cm}}
	\toprule
	$\HRm$ (Fraca) &
		Existe pelo menos um eixo de $W_i$ tal que ambos os extremos $E_1$ e
		$E_2$ são primos. O par $(E_1, E_2)$ é então um certificado de Goldbach
		para o par alvo $2M{+}2$. \\
	\midrule
	$\HRp$ (Forte) &
		Existe pelo menos um eixo de $W_i$ tal que a tripla completa
		$(E_1, \mathrm{Piv\hat{o}}, E_2)$ é inteiramente formada por primos.
		$\HRp$ implica $\HRm$. Esta condição mais forte conecta ao formalismo
		espectral do Artigo~6. \\
	\bottomrule
\end{tabular}
\end{center}

\medskip\noindent A separação entre camada algébrica e camada aritmética é
a essência da arquitetura do motor: a existência geométrica dos pares e sua
cobertura pelos eixos de $W_i$ são incondicionais; a primalidade simultânea
dos extremos é o único ingrediente aritmético. $\HRm$ é a lacuna central —
objeto do próximo artigo da série.

% ═══════════════════════════════════════════════════════════════════════════════
\section{O Gabarito — \texorpdfstring{$W_i$}{Wi} Canônica como Referência}
\label{sec:gabarito}

\subsection{O Que é o Gabarito}
\label{subsec:gabarito-def}

O Gabarito é uma instância canônica de~$W_i$: uma janela construída com os
pares obtidos pelo scanner na configuração canônica (ordenada) da grade,
contendo pelo menos um par primo nos seus eixos. Ele funciona como uma
fotografia da estrutura diagonal da fita registrada pelo caminho zigzag. É
um objeto fixo — referência geométrica — e não um estado dinâmico.

\begin{remark}[Analogia]
	Se o motor fosse um acelerador de partículas, o gabarito seria a placa
	de detecção com os padrões de colisão esperados pré-impressos. Quando os
	dados dinâmicos se alinham com o gabarito, o sensor dispara $\HRm$ ou
	$\HRp$.
\end{remark}

\subsection{Por Que o Gabarito Resolve o Problema de Estudo}
\label{subsec:gabarito-papel}

Quando $W_i$ é estudada diretamente na grade com dados aleatórios, os
eventos variam conforme a ordenação avança, tornando difícil isolar o
fenômeno de interesse. O gabarito resolve isso: mantém a estrutura das
diagonais fixa como referência, e durante a simulação compara-se o arranjo
dinâmico contra essa referência.

\begin{center}
	\fbox{\parbox{0.85\textwidth}{%
		\textbf{Estratégia experimental:}
		\begin{enumerate}[label=(\arabic*), leftmargin=1.5em, itemsep=0pt]
			\item Construir o Gabarito — $W_i$ canônica com pares obtidos pelo
			scanner na grade ordenada, com pelo menos um par primo.
			\item Iniciar o sistema em entropia máxima (caos puro).
			\item Aplicar ordenação gradual.
			\item A cada passo, sobrepor a configuração dinâmica ao Gabarito:
			quando as posições dos eixos coincidem com primos, o sensor dispara
			$\HRm$ ou $\HRp$.
		\end{enumerate}
	}}
\end{center}

% ═══════════════════════════════════════════════════════════════════════════════
\section{Protocolo de Simulação — Transição Caos $\to$ Ordem}
\label{sec:simulacao}

\subsection{Motivação}
\label{subsec:sim-motivacao}

O experimento de transição de fase serve a dois propósitos: (1) demonstrar
computacionalmente que $\HRm$ e $\HRp$ são ativadas durante a ordenação;
(2) medir o Tempo de Primeira Passagem — quantos passos (swaps) até que a
primeira janela ative o sensor.

\subsection{Protocolo}
\label{subsec:sim-protocolo}

\begin{center}
	\fbox{\parbox{0.85\textwidth}{%
		\begin{enumerate}[label=\arabic*., leftmargin=1.5em, itemsep=2pt]
			\item Escolher um par base $2M$ e construir a Fita-Dobra $2{\times}N$.
			\item Executar o scanner na grade canônica e construir o Gabarito
			$W_i$ (referência fixa).
			\item Embaralhar a lista de ímpares: entropia máxima (Caos Puro).
			\item Verificar no estado inicial: o caos puro já ativa o sensor por
			acaso?
			\item Aplicar ordenação total gradual (Bubble Sort — granularidade
			máxima).
			\item A cada swap: executar o scanner e comparar com o Gabarito.
			Registrar ativação de $\HRm$ e $\HRp$.
			\item Registrar: Tempo de Primeira Passagem e percentual do percurso
			total até a ativação.
		\end{enumerate}
	}}
\end{center}

\subsection{Métricas e Observações do Notebook}
\label{subsec:sim-metricas}

O notebook de verificação (\texttt{motor\_heranca.ipynb}) executa este
protocolo e coleta as seguintes métricas para cada par base testado:

\renewcommand{\arraystretch}{1.5}
\begin{center}
\begin{tabular}{>{\bfseries}p{4cm} p{9.2cm}}
	\toprule
	\textbf{Métrica} & \textbf{Descrição} \\
	\midrule
	Total de swaps & Número de trocas até a ordenação canônica completa. \\
	Passo $\HRm$ (swap) & Swap em que $\HRm$ é ativada pela 1ª vez — Tempo de Primeira Passagem. \\
	\% do percurso & Razão passo $\HRm$ / total de swaps: quão cedo no processo o evento ocorre. \\
	$\HRp$ ativada & Se existe $W_i$ com tripla prima completa $(E_1, \mathrm{Piv\hat{o}}, E_2)$ antes do fim. \\
	Nº de janelas $W_i$ & Janelas geradas pelo scanner $\approx \lceil(N{-}1)/4\rceil$. \\
	\bottomrule
\end{tabular}
\end{center}

\medskip\noindent Observações empíricas dos testes iniciais ($N$ pequeno):

\begin{center}
	\fbox{\parbox{0.85\textwidth}{%
		\begin{itemize}[itemsep=2pt]
			\item $\HRm$ é frequentemente ativada já no swap~0 — no estado de
			caos puro, antes de qualquer ordenação. A estrutura geométrica de
			$W_i$ captura pares primos mesmo com dados completamente embaralhados.
			\item Quando $\HRm$ não está ativa no caos, ela é ativada muito cedo
			no percurso (geralmente abaixo de 10\% do total de swaps),
			consistente com a Lei de Saturação Geométrica.
			\item $\HRp$ (tripla prima completa) é ativada em alguns casos, mas
			não em todos — esperado, pois é uma condição estritamente mais
			restritiva que $\HRm$.
			\item Os testes atuais cobrem $N$ pequeno. A verificação para $N$
			grande — com múltiplas janelas $W_i$ e análise do crescimento do
			Tempo de Primeira Passagem em função de $N$ — está planejada como
			próxima etapa computacional.
		\end{itemize}
	}}
\end{center}

% ═══════════════════════════════════════════════════════════════════════════════
\section{O Teorema do Motor (condicional a \texorpdfstring{$\HRm$}{HR-})}
\label{sec:teorema}

\begin{theorem}[Motor de Herança Estrutural]
	\textbf{Hipóteses:}
	\begin{itemize}[itemsep=2pt]
		\item[\textup{(HI)}] Existe $2M_0$ com par de primos $(p,q)$ tal que
		$p + q = 2M_0$ — par base certificado.
		\item[\textup{($\HRm$)}] Para toda iteração do motor, $W_i$ contém
		pelo menos um eixo com $E_1$ e $E_2$ ambos primos e
		$E_1 + E_2 = 2M_0{+}2$ — par alvo certificado por $\HRm$.
	\end{itemize}

	\textbf{Conclusão:} A cadeia
	\[
		2M_0 \;\to\; 2M_0{+}2 \;\to\; 2M_0{+}4 \;\to\; \cdots
	\]
	cobre todos os pares $\geq 2M_0$ como somas de dois primos. O par alvo
	de uma iteração torna-se o par base da iteração seguinte.
\end{theorem}

% ═══════════════════════════════════════════════════════════════════════════════
\section{Mapa de Resultados}
\label{sec:mapa}

\renewcommand{\arraystretch}{1.5}
\begin{center}
\begin{tabular}{p{5.8cm} >{\bfseries}p{2.8cm} p{4.8cm}}
	\toprule
	\textbf{Componente} & \textbf{Estatuto} & \textbf{Observação} \\
	\midrule
	Soma das colunas $= 2M$ (par base) &
		Incondicional &
		Álgebra pura; escolha de coluna prima certifica par base. \\
	Soma das diagonais $= 2M{+}2$ (par alvo) &
		Incondicional &
		$a_{(k+1)} + b_k = 2M{+}2$, $k \in \{1,\ldots,N{-}1\}$. \\
	Acoplamento $3C = 2N$ &
		Incondicional &
		Condição necessária e suficiente. \\
	$W_i$ fixa com 4 eixos de simetria &
		Incondicional &
		Estrutura geométrica pura. \\
	Scanner (pivôs borda$\to$centro, pula 1) &
		Incondicional &
		Preenche $W_i$ com diagonais do par alvo. \\
	Soma constante nos eixos de $W_i$ &
		Incondicional &
		$E_1 + E_2 = 2M{+}2$, herdado da fita. \\
	Gabarito como referência canônica &
		Incondicional &
		Foto da estrutura diagonal ordenada. \\
	$\HRm$ (par primo num eixo de $W_i$) &
		Condicional &
		Lacuna central — certifica o par alvo; objeto do próximo artigo. \\
	$\HRp$ (tripla prima completa em $W_i$) &
		Condicional &
		Versão forte; implica $\HRm$; conecta ao espectral. \\
	Teorema do Motor &
		Cond.\ a $\HRm$ &
		Cobre todos os pares $\geq 2M_0$ por indução. \\
	\bottomrule
\end{tabular}
\end{center}

\bigskip
\noindent\textit{Este documento é a base para a revisão dos Artigos~5 e~6
da série sobre a Conjectura de Goldbach. O código de simulação unificado
(\texttt{motor\_heranca.py}) implementa todos os módulos aqui descritos.}

\end{document}
